\documentclass[a4paper, 11pt]{article}


% =================================================================================
% ## MSP poster package
% =================================================================================
\usepackage{mspposter}

% =================================================================================
% ##
% =================================================================================
\usepackage[utf8]{inputenc}
\usepackage[T1]{fontenc}

% \usepackage[compact, center]{titlesec}
% \titleformat*{\section}{\large\bf\center}
% \titlespacing*{\section}{0pt}{0pt}{0.5\baselineskip}


% =================================================================================
% ##Title part
% =================================================================================

\usepackage{tabularx}
\usepackage{multicol}
\newcolumntype{C}[1]{>{\centering\arraybackslash}p{#1}}


\title{Example poster for the msp poster package}
\author{
    Alexander Hewer\textsuperscript{1-2}, Ingmar Steiner\textsuperscript{1-3}, Sébastien Le Maguer\textsuperscript{1,3}                                       \\
    \vspace{12pt}
    {\def\arraystretch{1.5}
    \begin{tabular}{C{0.35\textwidth} C{0.35\textwidth}}
      \textsuperscript{1}\textit{Cluster of Excellence MMCI, Saarland University} & \textsuperscript{2}\textit{DFKI Language Technology Lab, Saarbr\"{u}cken} \\
      \multicolumn{2}{c}{\textsuperscript{3}\textit{Computational Linguistics \& Phonetics, Saarland University}}                                                  \\
    \end{tabular}
    }
}
% =================================================================================
% ## Diverse commands
% =================================================================================

% reduce the vertical space generated by a section
\let \oldsection \section
\renewcommand{\section}{\vspace{-10pt}\oldsection}

\begin{document}

\begin{poster}

    % Set themes
    \setpostertheme{themes/default}

    \maketitle

    \begin{postercolumns}[2]

        \begin{posterblock}
            \begin{center}
                \begin{tikzpicture}
                    \node [fill=red!70, circle, text width=20pt] {};
                \end{tikzpicture}
            \end{center}
        \end{posterblock}

        \begin{posterblock}
            \begin{center}
                \begin{tikzpicture}
                    \node [fill=black!70, circle, text width=30pt] {};
                \end{tikzpicture}
            \end{center}
        \end{posterblock}

        \posternextcolumn

        \begin{posterblock}
            \begin{center}
                \begin{tikzpicture}
                    \node [fill=blue!70, circle, text width=40pt] {};
                \end{tikzpicture}
            \end{center}
        \end{posterblock}

        \begin{posterblock}
            \begin{center}
                \begin{tikzpicture}
                    \node [fill=orange!70, circle, text width=10pt] {};
                \end{tikzpicture}
            \end{center}
        \end{posterblock}

    \end{postercolumns}

    \begin{posterrows}[3]

        \begin{posterblockW}{150}{45}
            \section{Left Box}
            This poster block has a fixed height and width.
        \end{posterblockW}

        \begin{posterblockWFill}{35}
            \section{Right Box}
            Here, the poster block uses the remaining space that is available on the current row.
            Like the left block, it has a fixed height.
        \end{posterblockWFill}

        \begin{posterblock}
            {\bf posterblockWFill} automatically ends the row.
        \end{posterblock}

        \begin{posterblock}
            Blocks in this row are aligned to the largest block in the above row.
        \end{posterblock}

        \begin{posterblockH}{30}
            This poster block has a fixed height and uses the precomputed width.
        \end{posterblockH}

    \end{posterrows}

    \begin{postercolumns}[5]
        \begin{posterblock}
            \section{Left}
        \end{posterblock}
        \begin{posterblock}
            \section{Left}
        \end{posterblock}
        \begin{posterblock}
            \section{Left}
        \end{posterblock}

        \posternextcolumn
        \posternextcolumn

        \begin{posterblock}
            \section{Center}
        \end{posterblock}
        \begin{posterblock}
            \section{Center}
        \end{posterblock}
        \begin{posterblock}
            \section{Center}
        \end{posterblock}

        \posternextcolumn
        \posternextcolumn

        \begin{posterblock}
            \section{Right}
        \end{posterblock}
        \begin{posterblock}
            \section{Right}
        \end{posterblock}
        \begin{posterblock}
            \section{Right}
        \end{posterblock}

    \end{postercolumns}

    \begin{logobar}[logobar]
        \centering
        \fbox{\huge \bfseries LOGO 1}
        \hfill
        \fbox{\huge \bfseries LOGO 2}
        \hfill
        \fbox{\huge \bfseries LOGO 3}
        \hfill
        \fbox{\huge \bfseries LOGO 4}
    \end{logobar}

\end{poster}

\end{document}
